\chapter{\altGerEng{Zukünftige Arbeit}{Future Work}}

Many ways exist to improve this work and achieve a truly batteryless \ac{Li-Fi} receiver. Firstly, a more energy-efficient receiver circuit design is needed that performs well with a supply voltage lower than 2.6V. This could involve using lower-power electronics and designing \ac{EVM} operating on micro-level currents. Future prototypes can use a modulation scheme that works with smaller photodiode signals or be more adaptive to avoid saturation and noise. The power management unit needs a redesign, involving boost converters designed for low input power and the buck converter to handle the necessary power for receiver operation. Second, enhancing the \ac{EH} can be achieved by using photodetectors with better responsivity or larger area, or by using optical concentrators (lenses or reflectors) to focus more light onto the photodetectors.

Another fascinating direction would be to use time-based \ac{SLIPT}\cite{de2024reconfigurable}. For example, program the transmitter so that it occasionally sends a light pulse dedicated to charging, giving the harvester periodic boosts without interfering with data. Then, data transmissions can have their dedicated timing. This time-division approach could ensure the receiver gets enough energy without requiring a continuous high-power signal \cite{de2024reconfigurable} \cite{zhang2013mimo}. 

Fourth, exploring alternative light sources dedicated to improving \ac{EH} level\cite{de2024reconfigurable}. The thesis suggests using a laser or infrared LED for energy delivery in parallel with visible light communication. A laser can feed high-intensity light to the receiver’s photodetectors, increasing the \ac{EH} without affecting the data reception\cite{de2024reconfigurable}. This approach could overcome the limitations of LED light by employing a directional energy beam. 

Fifth, the implementation of the power-splitting prototype should be carried out further. This thesis work only provided theory and simulation. One could design a system where a photodiode or solar panel is biased to provide data AC coupling while still collecting DC energy.

Sixth, this thesis work only involved STOP Mode 0 while coding for STM32. There are STOP mode 1 and STOP mode 2, which consume less current: 6.70\(\mu\)A and 1.26\(\mu\)A. They can be integrated for more power efficiency\cite{pwr}.

Seventh, \ac{BER} analysis is necessary to better evaluate the communication branch \cite{fan2017circuit}\cite{Shin:16}\cite{de2024reconfigurable}. The \ac{BER} is determined by comparing the transmitted sequence of bits with the received bits and counting the number of errors \cite{fan2017circuit}. This ratio can be influenced by signal noise, distortion, and jitter\cite{ber}. Then, the Data rate can be measured against this \ac{BER} and harvested energy to plot using MATLAB, as done in \cite{liu2016visible}. It is also possible to compare the effect of distance on the system as shown in the paper \cite{liu2016visible}. In paper \cite{fan2017circuit}, the pseudo-random bit sequence generated by the transmitter FPGA and the comparator output are compared for calculating the \ac{BER} for different data rates.

Finally, further work should examine the system at larger scales and different scenarios: multi-LED networks, longer distances, and varying environmental lighting. Understanding how a network of \ac{Li-Fi} transmitters might keep a device charged would be valuable for practical deployment.


